\chapter{多模态理解}
\label{chap:vision-multimodal-understanding}

给一张街景照片，要求找出“右边拿杯子的人”，结果应落到一个可检查的区域。给一页票据，询问“总额是多少”，模型需要读出数字及其字段归属。给一段视频，追问“第二次敲门发生在讲话之前还是之后”，还需对齐事件、声音与时间。三个请求都可以返回一句话，但这句话背后所需的证据并不相同。

上一章已经解释了表示匹配、跨模态连接和语言解码。本章从请求要求的成果出发，说明这些接口怎样用于寻找候选、确定位置、恢复文字与结构，以及形成描述和答案。贯穿其中的是来源记录：某个判断依据哪幅图、哪一页、哪个区域或哪段时间。输入额外的OCR结果、字幕或真实数据表，会改变模型仍需完成的工作，因而也改变评价结论。

\section{跨模态匹配与检索}
\label{sec:vision-understanding-retrieval}

检索先把大量候选缩小到可以进一步观察的范围。查询可以是一句话，也可以是一幅图；候选可以是文字、图像或视频片段。模型提供比较分数之前，必须明确什么结果才算满足请求，否则高分只能说明某种相似，无法说明任务已经完成。

\subsection{查询与相关性的规定}
\label{sec:vision-understanding-relevance}

“穿红衣服的人正在骑车”同时规定了人物、衣服颜色和动作关系。画面中有红衣行人，也有另一人在骑车，只满足了部分条件；把红色归到自行车上，同样改变了请求。可把完整条件分解为对同一个对象$i$的判断
\[
  \operatorname{person}(i)\land\operatorname{redClothes}(i)
  \land\exists j\,[\operatorname{bicycle}(j)\land\operatorname{ride}(i,j)]\text{。}
\]
这个表达式是相关性协议的说明，不表示使用双编码器就会显式计算这些谓词。它的作用是防止标注和验收把“出现了所有关键词”当成“关系全部成立”。

候选单位也要固定。检索整段视频时，只要片中某处出现目标事件，就可能被标为相关；检索一个两秒片段，则要求该片段包含所需证据。一段文字可能描述多幅图片，一幅图片也可以有多种正确描述，因此应用数据的正例关系通常是多对多。训练划分还应避免同一视频的相邻片段或同一图的近重复版本分散到训练集和测试集。

为诊断组合关系，可以在保持对象词汇的条件下改变请求，如“人搬起箱子”和“箱子挡住人”，或把“戴帽子”改为“没有戴帽子”。相关性应按改变后的条件重新标注。部分相关候选可以保留等级，但不能在一个实验中按完整匹配计分，在另一个实验中又按同主题计分。

\subsection{候选内容的匹配}
\label{sec:vision-understanding-candidate-matching}

双编码路线把查询和候选分别变成表示，适合先缓存图库，再处理不断变化的请求。视觉输入怎样取样、哪种对应关系用于训练，会影响表示所保留的信息。以下方法都能形成排序，但其观察单位和配对依据不同。

\subsubsection{CLIP双编码的图文候选排序}
\label{sec:vision-understanding-clip}

沿用\ref{sec:vision-models-contrastive}节的归一化表示，把图库中第$i$幅图的向量保存为$\bm v_i$。文字查询$q$编码成$\bm u(q)$后，以
\[
  s_i=\bm u(q)^\top\bm v_i
\]
排序并返回图像编号。图像检索文字则缓存候选描述$\bm u_j$，编码查询图像$\bm v$，比较$\bm v^\top\bm u_j$。两种方向共用相似度形式，但候选集合、正例标注和评价分母不同。

假设查询向量为$(1,0)^\top$，三个图像向量为$(0.8,0.6)^\top$、$(0.6,0.8)^\top$和$(0,1)^\top$，均已归一化。教学排序为甲、乙、丙，分数依次为0.8、0.6、0。若人工核对发现乙才满足“红衣人物骑车”的全部关系，甲仅有相关对象，那么模型把部分相关候选排到了完整正例之前。将分数转成softmax并不会补上这个关系判断。

部署时要固定图像裁剪、文字模板、归一化方式和相似度口径。若用新的多对多相关性数据适配，可以令一个查询的全部已知正例参与目标，或者在采样时排除已知正例之间的负例冲突；具体目标属于新的适配设置。不能把原始CLIP的批次内单配对标注无条件当作图库的真实相关性。

\subsubsection{CLIP4Clip接口下的文字与视频语义匹配}
\label{sec:vision-understanding-clip4clip}

视频检索沿用\ref{sec:vision-video-clip4clip}节的帧编码、时间聚合和训练过程。使用时，把固定取样得到的视频表示与查询表示比较，输出视频编号及分数。图像已经提供的对象证据与视频增加的时间证据，应分别检查。

对同一图库，依次输入“箱子”“人搬起箱子”“穿红衣服的人搬起箱子”，请求增加了动作和参与者条件。排序改变可能来自真正的时间关系，也可能只来自某帧出现了红衣人物。可以定位命中片段，查看谁接触箱子、箱子是否离开支撑面，以及动作顺序。若帧聚合主要保留整段共同出现的内容，交换“人”和“箱子”的语义角色未必产生足够大的得分变化。

比较不同查询时保持抽帧、视频长度和候选范围不变。命中整段视频仅完成了候选选择，动作起止仍需时间定位。抽样漏掉搬起瞬间时，下游文字模型可能根据前后场景作合理猜测，但不能把该瞬间写成已经观察到的事实。

\subsubsection{HowTo100M的旁白配对与片段排序}
\label{sec:vision-understanding-howto}

教学视频的旁白提供了较易获得的图文联系。HowTo100M从自动转写的字幕行及其时间区间建立“视频片段—旁白文字”记录，再学习能够相互比较的表示。旁白可能提前说“接下来揉面”，也可能省略手上正在做的动作，因而这种记录是弱配对，不是逐帧动作标注\scite{vision-howto2019}

原实现使用预先提取并沿时间最大池化的视频特征，以及在词向量上计算一维卷积得到的文字特征。每一侧都经过线性变换与门控：若$\bm z=\bm W_1\bm x+\bm b_1$，则输出为$\bm z\odot\sigma(\bm W_2\bm z+\bm b_2)$，随后以两侧输出的余弦相似度$s_{i,j}$比较第$i$段视频和第$j$条旁白。门控按当前特征调节各维贡献，不需要两侧输入维数相同。

训练希望真实配对得分比错配至少高出间隔$\delta>0$。对一个配对$i$及负例$j$，两个方向的损失为
\begin{equation}
  L_{i,j}=[\delta+s_{i,j}-s_{i,i}]_+
          +[\delta+s_{j,i}-s_{i,i}]_+\text{，}
  \qquad [a]_+=\max(0,a)\text{。}
  \label{eq:vision-howto-ranking}
\end{equation}
第一项替换旁白，第二项替换视频。各项已经满足间隔时不再产生这项惩罚。原实现用批次内配对形成负例，并对同视频和跨视频贡献加权，使两类负例获得所设的平衡；这避免数量更多的跨视频负例压过同视频内的步骤区别。

以$\delta=0.1$、$s_{i,i}=0.7$、$s_{i,j}=0.65$、$s_{j,i}=0.4$为教学数值，损失为$0.05+0=0.05$。第一项要求把“同一厨房里的另一条旁白”区分开；只认出厨房背景还不够。使用时保存训练后的映射，对新旁白或查询编码，按片段分数排序，并把片段时间还原到原视频。排序目标鼓励区别配对，却不会自动消除旁白和实际动作之间的错位。

\subsection{精细对齐与结果排序}
\label{sec:vision-understanding-reranking}

独立编码便于搜索，细致关系判断则常需要让当前查询直接读取候选内容。可以把昂贵计算留给较小候选集，同时保留对召回错误的单独检查。

\subsubsection{BLIP-2的表示召回与配对重排}
\label{sec:vision-understanding-blip-retrieval}

BLIP-2的Q-Former在预训练时已经学过独立匹配和联合配对判定，机制见\ref{sec:vision-models-blip2}节。检索时，先取图像查询表示与文字表示之间的相似度，召回有限候选；再使候选文字与图像查询双向交互，由图文匹配头计算配对分数，重新排列这些候选。此路径不需要大型语言模型逐词生成答案。

原论文的检索适配共同微调图像编码器和Q-Former，并保留对比、匹配及图像条件文字生成三个目标；这不同于前述冻结图像编码器的预训练阶段。其推理设置先取128个候选，再按配对匹配分数重排。冻结范围和候选数都应随所用实验设置报告，不能从预训练说明中直接沿用。

重排能够让“同一画面里的红衣人物是否就是骑车者”影响联合计算，但仍需关系数据检验。若图库有三个完整正例，召回集仅含其中两个，那么无论重排多准确，完整正例的召回率最多为$2/3$。这个上限来自候选已经被删除，与后续模型有多少参数无关。

\subsubsection{局部对应与查询约束的复核}
\label{sec:vision-understanding-constraint-check}

对需要严格满足条件的查询，可以增加一个教学组合流程：从请求中列出需要核对的对象、属性、关系、计数或文字；在高分候选中调用定位和读取接口取得区域证据；再逐项记录满足、未满足或证据不足。组合规则不新增模型训练，也不属于CLIP原始输出。

例如，甲图得分0.8，但定位后发现红衣人物在步行，骑车者穿蓝衣，应记录动作绑定不成立。乙图得分0.6，相关人物和自行车位置、接触关系均可见，则可以作为更合适的结果。若车身遮住腿部，图像又无法区分骑行和推行，应保留歧义，而不是用较高相似度补足不可见部分。

结果可以同时返回候选编号、匹配分数、目标区域及满足的条件。图文排序、框或掩码定位、条件核验分别评价，才能知道错误产生在何处。下一节展开其中最直接的证据接口：把表达对应到具体位置。

\section{跨模态定位}
\label{sec:vision-understanding-grounding}

“右边拿杯子的人”需要落到画面中的对象，而非只生成一句改写。\term{视觉定位}（visual grounding）把文字或其他查询对应到视觉对象、空间区域或时间位置。输出的位置必须仍属于原始观测，才能用于复核、裁剪或后续操作。

\subsection{查询目标与参照条件的解析}
\label{sec:vision-understanding-grounding-query}

在上述表达中，“人”是目标类别，“拿杯子”是人与杯子的关系，“右边”还依赖参照系。若两个人都在拿杯子，“右边”可能指图像右侧，也可能指说话者右手方向。系统应把默认参照系写入协议；本章未另说明时，图像横轴向右、纵轴向下。

多图输入还需指定“图2中的人物”，页面输入需带页码，视频输入需带时刻或区间。“他旁边的杯子”依赖此前已经确定的“他”，不能仅按最近出现的名词替换。无法唯一确定时，可以返回候选集合或请求补充参照条件。查询目标根本不存在，也应允许输出空结果。

框、掩码、点击点和时间区间承担不同任务。框表示外接范围，掩码描述像素归属，点击点必须位于可操作控件内，区间描述事件持续时间。它们的几何关系可以辅助转换，但不能仅凭输出了四个数就认定完成框定位，也不能把框中心总是可点击当作普遍条件。

\subsection{查询与候选证据的对应}
\label{sec:vision-understanding-grounding-match}

定位需要保留空间或时间分辨率。把整图压成一个向量后再与整句比较，无法直接知道哪个位置支持哪段表达。下面分别考虑文字对应区域、声音对应区域，以及一段声音对应视频时间。

\subsubsection{MDETR的文字调制检测与短语对应}
\label{sec:vision-understanding-mdetr}

MDETR让图像特征与文字词元在检测过程中相互影响。图像经卷积骨干形成空间特征，文字经语言编码器形成词元特征；两侧线性投影到同一维数后沿序列方向拼接，进入联合Transformer编码器。对象查询读取联合表示，解码出框及其对应的文字位置，而不是只预测固定类别名称\scite{vision-mdetr2021}

训练记录包含图像、自由文本以及短语与框的对应。集合匹配沿用\ref{sec:vision-recognition-detr}节的思路，在预测框与真实框之间寻找一对一分配，再计算框的$L_1$损失和广义IoU损失。文字对应则允许多对多：一个短语可以指两个人，同一个人也可能在句中被不同短语指代。

设一个已匹配对象对应的文字词元位置集合为$T_i^+$，预测位置分布为$p_{i,j}$。软词元目标把概率质量均匀分给这些位置，其交叉熵为
\begin{equation}
  L_{\mathrm{token},i}
  =-\frac{1}{|T_i^+|}\sum_{j\in T_i^+}\log p_{i,j}\text{。}
  \label{eq:vision-mdetr-token}
\end{equation}
例如“红衣／人物”占两个词元位置，目标分别为$1/2$，不是要求模型任选一个位置得到概率1。未匹配对象查询预测“无对象”。实际对应建立在分词后的位置上，不能用中文字符数代替模型词元数。

仅预测位置还没有直接约束表示距离。MDETR另外投影并归一化对象表示$\bm o_i$与词元表示$\bm t_j$，计算对象到文字的对比项
\[
  L_{\mathrm{o\to t},i}
  =-\frac{1}{|T_i^+|}\sum_{j\in T_i^+}
    \log\frac{\exp(\bm o_i^\top\bm t_j/\tau)}
    {\sum_k\exp(\bm o_i^\top\bm t_k/\tau)}\text{。}
\]
文字到对象方向按相同方式对齐，再取两个方向的平均；没有正对应的项不参与该平均项的构造。总训练目标以权重组合框、软词元与表示对齐损失，共同更新相关编码和检测参数。软词元预测约束输出位置，对比项约束所用表示，二者使用相同标注但作用不同。

若图中有红衣人物甲和蓝衣人物乙，输入“红衣人物拿着杯子，蓝衣人物站在旁边”，模型可以同时输出两个人物框及对应短语位置。使用时按目标短语聚合有关位置得分，筛选框并恢复原图坐标；若任务只要求拿杯者，还需排除作为参照出现的乙。框能支持区域检查，像素掩码需要另加分割路径，复杂关系能否正确消歧仍取决于训练与测试覆盖。

\subsubsection{Grounding DINO的查询组织与短语定位}
\label{sec:vision-understanding-grounding-dino}

Grounding DINO的跨模态特征、查询选择和解码已经在\ref{sec:vision-recognition-grounding-dino}节展开。用于定位时，把“人物”“戴帽子的人”“右侧戴帽子的人”等不同范围的表达作为文字输入，得到预测框与文字位置的匹配分数。类别词列表应按所用接口的分隔约定组织，完整指代表达则应保留属性和关系，不随意拆成互不相关的词。

一种常见使用过程是以框内最高词元分数和框阈值筛出候选，再以文字阈值选择相关词元并恢复短语。两道阈值作用不同：降低框阈值会增加候选区域，降低文字阈值会扩大附着在框上的文字范围。二者都不能把“右侧”一定指向当前目标作为逻辑保证，返回的短语也未必是连续、完整的语法单位。

对“右侧戴帽子的人”，应核对帽子属于哪个人物、“右侧”的参照是否一致，以及是否存在第二个满足条件的候选。需要概念对应的实例掩码时，可以使用\ref{sec:vision-recognition-sam3}节的SAM 3接口，分别检查概念存在与实例范围。框和掩码的组合增加了成果类型，并没有免除关系核验。

\subsubsection{困难区域对比的可见声源定位}
\label{sec:vision-understanding-sound-source}

一帧画面中同时出现说话者和旁观者，整图与语音相似还不足以区分声源。Localizing Visual Sounds the Hard Way研究了这一问题\scite{vision-visual-sounds2021}
该方法保留图像特征的空间网格，把音频频谱编码为向量，再逐位置比较，形成可见声源热图。

记图像位置$(u,v)$的非零特征为$\bm f_{u,v}\in\mathbb R^d$，非零音频特征为$\bm a\in\mathbb R^d$，对应分数为
\begin{equation}
  S_{u,v}=\frac{\bm a^\top\bm f_{u,v}}
                 {\lVert\bm a\rVert_2\lVert\bm f_{u,v}\rVert_2}\text{。}
  \label{eq:vision-sound-map}
\end{equation}
使用时，将这张网格热图映射回图像，并按所定阈值与可见声源区域比较。声音编码器输出整段表示，因此一张热图描述的是这段声音与当前图像的空间联系，不是每个音频采样点的位置。

训练只有音画配对时，真实声源掩码不可用。方法用两道平滑阈值构造较可信的正区域和同图负区域：
\[
  m^+_{u,v}=\sigma\!\left(\frac{S_{u,v}-\epsilon_+}{\tau}\right),
  \qquad
  m^-_{u,v}=1-\sigma\!\left(\frac{S_{u,v}-\epsilon_-}{\tau}\right),
  \qquad \epsilon_+>\epsilon_-\text{。}
\]
高分位置获得较大正权重，低分位置获得较大负权重，中间区域对两侧的贡献较小。由于使用sigmoid，中间带不是严格全部为零。正响应是以$m^+$加权的归一化平均相似度，负响应同时包含以$m^-$加权的同图背景和其他视频的错配响应；对比目标提高正响应相对于负响应的得分，更新图像和音频编码路径。

同图困难负例使训练需要区分发声者和视觉上同样显著的旁观者，而不只区分“有人说话的画面”与“另一种场景”。这种自举仍可能把错误的高分区域当成正例。若声音来自画外，仅输出热图峰值就会强行指定一个可见对象；处理该情况需要另设“无可见声源”判定和相应数据，不能从原热图接口中直接得到保证。

\subsubsection{音画距离学习与跨模态时间定位}
\label{sec:vision-understanding-av-time}

空间声源回答“画面哪里”，跨模态时间定位回答“另一模态的哪一段与这段内容同步”。Audio-Visual Event Localization工作为此使用音画距离网络：分别提取一秒音频与视觉片段特征，经各自两层全连接投影到共同维数，以欧氏距离$D$比较\scite{vision-ave2018}

若$y=1$表示同步配对，$y=0$表示错配，训练目标为
\begin{equation}
  L=yD^2+(1-y)[m-D]_+^2\text{，}\qquad m>0\text{。}
  \label{eq:vision-av-distance}
\end{equation}
同步对被拉近，错配对在距离小于$m$时被推开。标签来自训练配对关系，不是模型使用时预先知道目标时间。

设音频查询含$l$个时间单元，候选视频含$T$个时间单元。沿视频滑动同样长度的窗口，选择累计距离最小的起点
\begin{equation}
  \hat r=\argmin_{1\leq r\leq T-l+1}
    \sum_{j=1}^{l}D(\bm v_{r+j-1},\bm a_j)\text{。}
  \label{eq:vision-av-sliding}
\end{equation}
这里$\bm v_k$、$\bm a_j$表示已投影的片段特征。若单元时长为$\Delta$，且视频从0秒开始，第$r$个单元覆盖$[(r-1)\Delta,r\Delta)$，输出区间为$[(\hat r-1)\Delta,(\hat r+l-1)\Delta)$。

图\ref{fig:vision-av-temporal-match}用二维投影的教学数据走通三次窗口比较。取音频向量$(0,0)$、$(1,0)$，四段视频向量$(2,0)$、$(0,0.1)$、$(1,0.2)$、$(3,0)$。三个窗口的累计距离约为3.005、0.300和3.020，故选中第2、3单元；当$\Delta=1$秒时，区间为$[1,3)$秒。

\begin{figure}[htbp]
  \centering
  % 教学欧氏距离：声音(0,0),(1,0)，视频(2,0),(0,.1),(1,.2),(3,0)。
\begin{tikzpicture}[x=1cm,y=1cm,font=\figurefont\small]
  \node[anchor=west] at (0,3.5) {逐单元距离$D(\bm v_k,\bm a_j)$};
  \foreach \x/\k in {2.2/1,4.2/2,6.2/3,8.2/4}
    \node at (\x,2.7) {视频$\k$};
  \node[anchor=east] at (1.1,1.8) {声音1};
  \node[anchor=east] at (1.1,0.7) {声音2};
  \foreach \x/\y/\v in {
    2.2/1.8/2.000,4.2/1.8/0.100,6.2/1.8/1.020,8.2/1.8/3.000,
    2.2/0.7/1.000,4.2/0.7/1.005,6.2/0.7/0.200,8.2/0.7/2.000}{
    \draw[BookRule,line width=0.5pt,fill=BookPaper]
      (\x-0.9,\y-0.4) rectangle (\x+0.9,\y+0.4);
    \node at (\x,\y) {$\v$};
  }
  \draw[BookTeal,line width=1.2pt] (3.3,1.4) rectangle (5.1,2.2);
  \draw[BookTeal,line width=1.2pt] (5.3,0.3) rectangle (7.1,1.1);
  \node[anchor=west] at (0,-0.2) {窗口累计距离：};
  \node[align=center] at (1.7,-1.0) {视频1--2\\$3.005$};
  \node[align=center,text=BookTeal] at (4.8,-1.0) {视频2--3\\$\mathbf{0.300}$};
  \node[align=center] at (7.9,-1.0) {视频3--4\\$3.020$};
  \node[anchor=west] at (0,-2.0) {选中$[1,3)$秒；时间单元长1秒};
\end{tikzpicture}

  \caption{音频查询与视频窗口对应的教学计算。每列是一个视频时间单元；窗口内按相对次序逐项相加，粗框标出累计距离最小的窗口。向量和距离均为构造值，不是模型测量结果。}
  \label{fig:vision-av-temporal-match}
\end{figure}
\FloatBarrier

重复敲门可能在两处形成同样小的距离，最小值不再唯一；较粗的单元又限制边界精度。原工作另外研究逐秒预测“同时可见且可听的事件类别”，其监督是事件标签，输出是时间标签序列。该任务与上述查询对应的监督及输出均不同，不能把两套结果相互替代。

\subsection{关系、状态与上下文消歧}
\label{sec:vision-understanding-context}

截图中可能有两个“保存”按钮，一个属于当前对话框，另一个属于背景窗口。文字与局部外观都相似时，必须把窗口归属和当前状态一起纳入定位。高分辨率还带来另一限制：整屏缩小后，决定控件区别的小字可能消失。

\Needspace{60mm}
\subsubsection{SeeClick的截图指令与坐标预测}
\label{sec:vision-understanding-seeclick}

SeeClick以截图、控件描述与位置构成训练记录，把坐标作为可生成的文字\scite{vision-seeclick2024}
例如，输出可以写成\code{click (0.49, 0.40)}，两个数表示相对于截图宽、高的归一化位置。其原始实现以Qwen-VL为基础，用LoRA适配视觉编码器和语言模型，并混合网页、移动界面及一般视觉语言数据。

监督仍是目标输出序列的交叉熵，模型逐词预测位置文本，不是增加一个欧氏距离回归头。数字生成损失并不直接等于像素误差，因此使用时还要解析数值、检查范围和坐标约定。若截图宽1600、高1000像素，输出$(0.49,0.40)$对应连续坐标$(784,400)$；实际点击还需与显示缩放和截图坐标一致。

定位评价可检查点是否落在目标控件框内。点落在正确按钮上，并不能说明该按钮当前可用，更不能说明保存已经成功。多步操作还需在新截图或状态中确认结果。这里的任务只对当前截图产生定位，历史动作和界面状态变化应作为下一次观测重新进入系统。

\subsubsection{ScreenSeekeR的规划裁剪与逐级搜索}
\label{sec:vision-understanding-screenseeker}

将高分辨率整屏压缩到固定输入尺寸时，小控件可能只剩几个像素。ScreenSeekeR利用界面的层级结构缩小搜索区域：规划模型根据指令提出目标或邻近控件可能所在的位置，定位模型将这些描述落到框，候选区域经扩张、评分与去重后进入递归搜索。ScreenSpot-Pro是评测基准，ScreenSeekeR是论文提供的搜索方法，二者名称不能互换\scite{vision-screenseeker2025}

候选区域的评分不是直接采用规划模型的口头置信度。原方法让定位框的中心为候选区域投票：将中心$(x,y)$转换为区域内部归一化位置$(x',y')$，在区域内时给予
\[
  s=\exp\!\left(-\frac{(x'-0.5)^2+(y'-0.5)^2}{2\sigma^2}\right)
\]
的贡献，区域外贡献为0，原设置取$\sigma=0.3$。将所有定位框的贡献相加，再用非极大值抑制减少重复候选。靠近区域中心的多条线索可以提高该区域的优先级，但这种几何投票仍可能支持错误窗口。

搜索按评分进入裁剪区域，继续提出位置或进行直接定位，直到规划模型确认目标、区域达到所定尺寸条件或到达最大深度。每次裁剪都应保存到父图的坐标变换，再组合回全屏。例如，父裁剪左上角为$(400,200)$，其内再裁出左上角$(100,50)$、宽高$200\times100$的区域；若最终归一化点为$(0.25,0.60)$，全屏点为$(550,310)$。若中途调整过尺寸，还要先逆转该次缩放。

该方法组合既有模型，不增加训练。它把预算从一次全屏读取改为多次局部搜索，结果应同时报告调用次数、像素或词元预算和停止条件。裁剪虽然放大了目标，也可能删去“当前窗口”等消歧线索，因此递归状态仍需保留整屏上下文与候选所属位置。

\subsection{定位结果与歧义反馈}
\label{sec:vision-understanding-grounding-output}

一个可复核的定位结果至少包含查询、源图或页面编号、时间信息、坐标约定及输出范围。框可以继续作为分割提示，点可以交给界面执行器，但新的输出要按新的标准检查。返回多个合理候选时，应说明差别来自参照不明、重复实例，还是目标部分遮挡。

视频中“刚才拿杯子的人”还要求身份跨帧延续。可以复用\ref{sec:vision-video-tracking}节的对应和存在状态，区分进入、持续可见与离开；离开画面的对象不应被强行重新绑定到一个相似人物。2026年公开的TempoGround将跨帧对应、对象存在及二维和相机坐标系中的三维框联系起来，可作为这一问题的扩展案例。其二维定位、身份保持和三维几何仍是分别受约束的结果，论文设置不构成任意视频的长期一致性保证\scite{vision-tempoground2026}

\section{文字读取与结构化解析}
\label{sec:vision-understanding-reading}

找到票据区域以后，仍需区分“看见数字”和“知道数字属于哪个字段”。表格中，即使所有字符都读对，合并表头归属错误仍会改变含义。读取和解析因此需要逐步恢复字符串、页面关系、单元结构和连接，而不是只输出一段外观相似的文字。

\subsection{字符与字符串的恢复}
\label{sec:vision-understanding-ocr}

\term{光学字符识别}（optical character recognition, OCR）把文字图像转换为机器可处理的字符串。完整流程通常还需文字检测，即找出哪些区域包含文字。检测结果为框或多边形，识别结果为字符串；位置正确而内容读错、内容正确但属于另一行，是两种不同错误。

\subsubsection{CRAFT的字符区域与关联分组}
\label{sec:vision-understanding-craft}

弯曲招牌中的字符未必落在同一条直线上，直接拟合一个文字矩形可能包含大量背景。CRAFT预测两张稠密图：字符区域图给出各字符中心附近的响应，关联图给出相邻字符之间应当连接的位置。网络以卷积骨干和多层特征融合恢复空间分辨率，末端两个通道分别承担这两种预测\scite{vision-craft2019}

合成数据有字符框，可把二维高斯模板通过透视变换映射到字符框内，形成区域目标；相邻字符框之间的区域形成关联目标。真实数据常只有文字级框和字符串，原方法先裁出文字，用当前模型的区域响应和分水岭划分估计字符框，再把坐标还原到原图，作为弱监督。

这种估计并非无误。原方法利用已知字符串的长度与预测字符数之差衡量伪标注可靠性，对区域和关联图的平方误差加权；过低可靠性的预测改用按字符数均分文字区域的退化估计。合成字符监督和带文字级标注的真实样本共同参与训练，使“从字符连成文字”的计算有可学习目标。

推理时，将超过区域阈值或关联阈值的位置并入二值图，执行连通分组，再为各组恢复外接四边形或适合弯曲文字的多边形。教学招牌“青山路”若三个字符响应之间有两条正确关联，会形成一个文字实例；若“路”与旁边门牌之间也出现强关联，分组就可能错误合并。此时问题发生在文字定位，后续识别器再准确也未必能恢复正确实例边界。

\subsubsection{TrOCR的文字行编码与词元解码}
\label{sec:vision-understanding-trocr}

文字位置确定后，TrOCR把裁剪文字图像转成字符串。图像编码器对图像块计算表示，文字解码器读取这些表示与已有词元，预测后续词元。原论文的典型配置把文字图像调整到$384\times384$，使用$16\times16$图像块；这一固定尺寸设置与后文Pix2Struct的可变分辨率设计不同\scite{vision-trocr2022}

编码器与解码器分别利用已有视觉和语言参数初始化，新增的跨注意力连接也要通过文字识别数据学习。以文字图像$\bm X_{\mathrm{text}}$和目标词元$y_1,\ldots,y_m$为一条记录，训练最小化
\begin{equation}
  L_{\mathrm{OCR}}
  =-\sum_{j=1}^{m}\log p_{\bm\theta}
    (y_j\mid y_{<j},\bm X_{\mathrm{text}})\text{。}
  \label{eq:vision-trocr-loss}
\end{equation}
训练包含合成文字及目标任务适配。使用时从起始符逐词解码，直到结束符或长度上限；可以用束搜索保留若干候选，最终去掉特殊词元并恢复字符串。这里的预测单元可以是子词，不要求每一步恰好一个字符。

模糊字形“1”和“I”可能在视觉上相近，语言上下文会影响选择。对常见地名，这有时有助于消歧；对票据编号或未见专名，强语言先验也可能把正确的罕见字符串改成常见词。应保留原裁剪和识别结果用于核对。TrOCR处理文字图像，并不自行给出全页文字位置、双栏阅读顺序或表格归属。

\subsection{页面区域与阅读关系的恢复}
\label{sec:vision-understanding-page}

一页文档上的文字可能属于标题、正文、表格、脚注或公式。把左栏第一行和右栏第一行交错拼接，即使每个字都正确，仍会得到错误句子。页面编码需要利用二维位置，结构输出还要明确定义元素及阅读关系。

\subsubsection{LayoutLMv3的文字、布局与页面编码}
\label{sec:vision-understanding-layoutlm}

LayoutLMv3将OCR词元、对应二维框以及页面图像块共同送入Transformer。文字嵌入包含内容、序列位置和布局位置，图像块经线性投影并加入位置表示。两类单元在共同上下文中交换信息，使“128.50”可以同时参考附近的“总额”及其页面位置\scite{vision-layoutlmv32022}

预训练包含三种互补目标。文字掩码目标预测被隐藏的文字词元；图像掩码目标预测被隐藏图像块的离散视觉编码，而不是直接回归原始像素；文字—图像块对齐目标则在未被遮盖的文字位置上，预测其对应图像块是否被遮盖。后一个目标利用OCR位置提供的联系，不是判断任意一段文字与整幅图是否语义一致。

若一行金额文字位于图像块7和8，遮盖相应块而保留词元时，对齐标签随之改变。这样的监督促使表示利用内容与局部图像的对应，但不能保证OCR本身已经读对。下游字段标注还需在文字表示上训练任务头，按每个词元或文字块预测“编号”“日期”“总额”等类别。

对于版面检测，原论文也使用不输入文字嵌入的视觉任务路径，以检测头输出标题、表格等区域。它与需要OCR的文字字段任务是不同配置。LayoutLMv3因此不能被理解为一个无条件输出整页HTML及阅读顺序的解码器；具体能恢复什么，仍由输入、任务头和监督决定。

\subsubsection{Pix2Struct的可变分辨率与结构序列生成}
\label{sec:vision-understanding-pix2struct}

直接把页面像素转换为结构文字，可以减少使用时对外部OCR的依赖，但必须同时保留小字与布局。Pix2Struct按输入词元预算缩放图像，尽量保持原纵横比，并选取预算内尽可能多的固定尺寸图像块，再用二维位置区分块的行列。图像编码器输出供文字解码器生成目标序列\scite{vision-pix2struct2023}

预训练从网页截图及HTML建立对应。原方法保留可见内容相关的结构，压缩冗长的嵌套，只保留所选节点的文字与图像信息，再选取能放入目标长度的子树。截图上标出该子树的区域，并遮盖其中部分文字，训练解码器恢复完整的简化结构。阅读热身阶段先学习从合成文字图像恢复字符串，再进入截图解析阶段。

因此，预训练目标是经过筛选和简化的结构串，不能直接等同于包含全部CSS、坐标和交互行为的原网页HTML。若教学页面具有“发票—编号A017—总额128.50”的层次，可以用嵌套节点表示从属关系；解析器应检查节点是否闭合、文字是否遗漏及金额属于哪个字段。字符串符合语法，只说明可以读入结构，不说明内容忠实于页面。

页面元素恢复后，还要按既定关系组织阅读顺序。双栏正文应在本栏顺读，再转入下一栏；跨栏标题和所属公式要与相应内容连接。OmniDocBench同时提供完整文档解析和文字、版面、表格、公式等模块的评价\scite{vision-omnidocbench2025}
这些分项使整页误差能够定位到具体结构，而非仅用一串文字相似度概括。按页面类型保留结果，有助于发现报纸密集跨栏和论文公式等不同困难。

用于下游问答时，Pix2Struct把问题等条件渲染到输入图像上，再以答案作为目标适配；界面任务也可以把目标区域等条件画入输入。使用时必须遵循相应检查点的预处理和输出约定。改变词元预算会改变可保留的视觉细节，解码更长的字符串不能恢复已经因缩小而不可辨认的小数点。

\subsection{表格单元与字段关联的恢复}
\label{sec:vision-understanding-tables}

表格中，一个数的含义来自行、列和表头。票据字段则来自键与值的对应。二者都需要把空间区域组织成有语义的关系，不能只把OCR字符串按出现顺序排列。

\subsubsection{PubTables-1M的表格检测与结构对象恢复}
\label{sec:vision-understanding-pubtables}

PubTables-1M将表格检测与内部结构恢复分别建模。一个DETR模型在页面定位表格，另一个在裁剪表格中预测整表、行、列、列表头、投影行表头和跨格单元等对象。投影行表头是横跨其他列、用于划分后续行组的标题。两个模型都使用检测监督，集合匹配原理沿用\ref{sec:vision-recognition-detr}节\scite{vision-pubtables2022}

行和列对象相交形成基本格子。跨格单元指出哪些基本格应合并，表头对象提供其功能范围。训练标注的规范化还处理过度切分，例如把本应覆盖两列的“金额”误标成一格文字和一格空白。规范化约定让同一种表头含义对应稳定结构，但不能保证所有文档都遵守相同排版习惯。

图\ref{fig:vision-table-structure}用四行三列的教学表说明恢复过程。前两行为表头：“项目”跨两行，“金额（元）”跨第2、3列，其下分别为1月、2月。后两行是甲的$(12,18)$与乙的$(8,10)$。行列交叉先产生12个基本格，两次合并后形成10个逻辑单元，跨格表头使每个金额都能回溯到月份和单位。

\begin{figure*}[htbp]
  \centering
  % 构造表：4行3列；两个跨格对象各合并2格，得到10个逻辑单元。
\begin{tikzpicture}[x=1cm,y=1cm,font=\figurefont\small,
  arr/.style={-{Stealth[length=2mm]},BookInk,line width=0.9pt}]
  \node[anchor=west,font=\figurefont\bfseries\small] at (0,4.0)
    {行列交叉与跨格对象};
  \node[anchor=west,font=\figurefont\bfseries\small] at (9,4.0)
    {恢复后的逻辑表};
  \foreach \dx in {0,9}{
    \begin{scope}[xshift=\dx cm]
      \fill[BookBlue!12] (0,1.6) rectangle (2,3.2);
      \fill[BookTeal!12] (2,2.4) rectangle (6,3.2);
      \draw[BookInk,line width=0.9pt] (0,0) rectangle (6,3.2);
      \draw[BookRule,line width=0.5pt] (0,0.8) -- (6,0.8);
      \draw[BookRule,line width=0.5pt] (0,1.6) -- (6,1.6);
      \draw[BookRule,line width=0.5pt] (2,0) -- (2,3.2);
      \node[fill=BookBlue!12,inner sep=2pt] at (1,2.4) {项目};
      \node[fill=BookTeal!12,inner sep=2pt] at (4,2.8) {金额（元）};
      \foreach \x/\y/\v in {
        3/2/1月,5/2/2月,1/1.2/甲,3/1.2/12,5/1.2/18,
        1/.4/乙,3/.4/8,5/.4/10}
        \node at (\x,\y) {\v};
    \end{scope}
  }
  \draw[BookMuted,dashed,line width=0.5pt] (4,0) -- (4,3.2);
  \draw[BookMuted,dashed,line width=0.5pt] (0,2.4) -- (6,2.4);
  \draw[BookBlue,line width=1.1pt] (0,1.6) rectangle (2,3.2);
  \draw[BookTeal,line width=1.1pt] (2,2.4) rectangle (6,3.2);
  \node[fill=BookBlue!12,inner sep=2pt] at (1,2.4) {项目};
  \node[fill=BookTeal!12,inner sep=2pt] at (4,2.8) {金额（元）};
  \draw[BookRule,line width=0.5pt] (13,0) -- (13,2.4);
  \draw[BookRule,line width=0.5pt] (11,2.4) -- (15,2.4);
  \draw[arr] (6.3,1.6) -- (8.6,1.6);
  \node[align=center] at (7.45,2.35) {按跨格对象\\合并};
  \node at (3,-0.6) {12个基本格};
  \node at (12,-0.6) {10个逻辑单元};
\end{tikzpicture}

  \caption{行列检测到逻辑单元的教学恢复。左侧虚线表示行列交叉形成的基本格；右侧根据两处跨格对象合并表头。表头与单元归属是结构成果，字符内容来自另行读取；图中的数值供后续问答复算。}
  \label{fig:vision-table-structure}
\end{figure*}
\FloatBarrier

推理还要解决结构冲突，例如两个跨格对象同时声称拥有同一基本格。原实现通过抑制冲突对象或调整范围，使行、列及跨格关系可以组成逻辑表；之后将有位置的文字分配给单元，输出行列索引、跨度、内容和来源框。不能只输出检测框就声称已经恢复HTML表格。

表格结构正确与字符串正确应分别评价。若甲2月的“18”被读成“13”，结构可以完全正确而后续差值错误；若“金额”被当作只覆盖1月，所有字符正确也仍损失单位归属。原工作不覆盖跨页表格，连续两页的表头重复和行连接需要额外处理。

\subsubsection{布局字段标注与键值对应}
\label{sec:vision-understanding-fields}

票据的“编号”“日期”“总额”不是规则二维数据表，可以复用LayoutLMv3页面表示预测文字块角色，再按内容、位置和字段约束连接键与值。角色标注需要对应的字段监督；连接可以由另训关系头完成，也可以使用清楚写出的规则，不能把后者无条件归为LayoutLMv3原有结构。

设教学票据的读取结果含“编号A017”“日期2026-04-03”“已付100.00”“总额128.50”。一种简单规则是在每个键的同行右侧寻找值，要求日期符合日期格式、金额具有数值格式，并避免同一值被互不相容的字段重复占用。这条规则可以在该版式下恢复记录，但遇到“值在键下方”的模板就需调整。

字段结果应保留原文字、规范化值和区域来源。例如，原文“128.50”对应数值128.50及单位元；不能因另有“已付100.00”就将两者都归为总额。实体与关系抽取的通用方法可参见\BookPartRef{part:natural-language-processing}{“自然语言处理”部分}，这里增加的是字段如何落在页面证据上。

\subsection{图表数据与编码关系的恢复}
\label{sec:vision-understanding-charts}

图表用位置、长度、颜色或角度表达数值。读取它需要反向还原这些编码，再把数值对应到标签和系列。只恢复一个高度而不知道它属于哪个月份，仍不能回答月份比较。

\subsubsection{ChartQA中的关键点读数与标签关联}
\label{sec:vision-understanding-chart-reading}

ChartQA采用的ChartOCR式恢复管线先检测绘图区、条形端点、折线顶点及文字区域，再结合文字识别估计刻度，并用位置和颜色关联标签与图例\scite{vision-chartqa2022}
条形可由左上和右下关键点恢复，折线由连接点恢复；不同图元需不同的几何读法。

在线性纵轴上，已知像素纵坐标$y_0,y_1$分别对应数值$v_0,v_1$，则位置$y$的读数为
\begin{equation}
  v(y)=v_0+\frac{y-y_0}{y_1-y_0}(v_1-v_0)\text{，}
  \qquad y_1\neq y_0\text{。}
  \label{eq:vision-chart-calibration}
\end{equation}
这里纵坐标向下增加，因此较大的数值通常对应较小的$y$。公式来自两个锚点间的线性插值，不能用于未变换的对数轴。

沿用表格中甲的金额，设图表0元刻度位于$y_0=340$，20元刻度位于$y_1=140$，两根柱顶分别为220和160像素。代入得到12元和18元。检测结果还需依据横轴位置对应1月、2月，再依据图例把它们归到甲系列。图\ref{fig:vision-chart-reading}突出这个“像素—刻度—记录”的联系。

\begin{figure}[htbp]
  \centering
  % 数值与table-structure-recovery相同；像素轴 y=340-10v。
\begin{tikzpicture}[x=1cm,y=1cm,font=\figurefont\small,
  arr/.style={-{Stealth[length=2mm]},BookInk,line width=0.8pt}]
  \node[anchor=west] at (0,4.8) {金额（元）};
  \draw[arr] (0.8,0) -- (0.8,4.4);
  \draw[arr] (0.8,0) -- (5.0,0);
  \foreach \y/\v in {0/0,2/10,4/20}{
    \draw[BookMuted,line width=0.5pt] (0.7,\y) -- (0.8,\y);
    \node[anchor=east] at (0.65,\y) {\v};
    \draw[BookRule,dashed,line width=0.5pt] (0.8,\y) -- (4.7,\y);
  }
  \draw[BookBlue,line width=0.9pt,fill=BookBlue!18]
    (1.65,0) rectangle (2.35,2.4);
  \draw[BookBlue,line width=0.9pt,fill=BookBlue!18]
    (3.35,0) rectangle (4.05,3.6);
  \node[anchor=south] at (2,2.5) {$y=220$};
  \node[anchor=south] at (3.7,3.7) {$y=160$};
  \node[anchor=north] at (2,-0.1) {1月};
  \node[anchor=north] at (3.7,-0.1) {2月};
  \node[anchor=north] at (2.9,-0.7) {甲系列};
  \node[anchor=west,font=\figurefont\bfseries\small] at (5.5,3.7)
    {恢复记录};
  \node[anchor=west] at (5.5,2.7) {甲，1月，12元};
  \node[anchor=west] at (5.5,1.7) {甲，2月，18元};
  \node[anchor=west,align=left] at (5.5,0.5)
    {像素纵坐标向下\\金额沿纵轴向上};
\end{tikzpicture}

  \caption{线性纵轴的教学读数。轴上标注金额，柱顶旁标注原图像素纵坐标；右侧记录同时保留系列、月份、数值与单位。两根柱的数值与图\ref{fig:vision-table-structure}中甲行一致。}
  \label{fig:vision-chart-reading}
\end{figure}
\FloatBarrier

若纵轴是对数刻度，应在线性关系成立的$\log v$上插值，再取指数恢复$v$；若柱从非零基线开始，只比较柱长又会改变读数。颜色近似、图例遮挡和折线交叉会造成系列归属错误，需要独立于数值误差检查。

原ChartQA比较真实数据表可用与从图表恢复数据表两种设置。前者提供了额外的数值和结构信息，后者把读数误差传到问答。将“18”误读成“13”后，计算器仍能准确算出$13-12=1$，却无法得到真实差值6；算术正确不能修复此前丢失的视觉证据。

\subsection{图解节点与连接关系的恢复}
\label{sec:vision-understanding-diagrams}

流程图中的箭头表示有向关系，位置相邻的两个框未必相连。图解解析需要恢复节点及其文字，还要追踪连接端点和方向，才能形成后续推理可用的图。

\subsubsection{节点检测、连线追踪与关系图构造}
\label{sec:vision-understanding-diagram-graph}

一种教学组合方法先用区域检测取得节点框，在各框内识别文字；随后检测线段和箭头尖端，沿连续线迹追踪连接；再依据端点与节点边界的接触，将线迹归到源节点和目标节点。节点检测可用框监督，连接预测可用端点或关系监督；在规则线稿中，也可以明确定义阈值、线段连接和箭头形状规则而不训练模型。

设节点集合为$V=\{A,B,C,D\}$，目标输出为有向图$\mathcal G=(V,E)$。若两条连线$A\to D$和$B\to C$在中间交叉，交点处没有连接点标记，恢复结果仍应只有这两条边；按“距离最近”接线可能误造$A\to C$和$B\to D$。断线、遮盖或缺少箭头时，应保留方向未知或多个候选，不能仅由页面上下顺序决定方向。

输出记录可以写成“源节点A，边标签通过，目标节点D”，并带原始线迹和端点位置。结构检查包括是否遗漏边、方向是否反转、边标签是否属于当前连接。这个流程是检测、识别与几何规则的组合，使用Pix2Struct生成某种结构串时也仍需相应监督和验证，不能将其视为原论文固定提供的通用图解析器。

\section{视觉描述与内容概括}
\label{sec:vision-understanding-description}

定位和读取保留具体证据，描述则选择其中一部分组成语言。描述可以覆盖全图、指定区域或一段事件，覆盖范围决定应当观察哪些内容。语言解码仍沿用\ref{sec:vision-models-language-generation}节的条件生成，当前需要补足的是内容如何被选择，以及选择后如何保持与来源一致。

\subsection{场景与全局内容的概括}
\label{sec:vision-understanding-global-caption}

街景短描述可能只需指出“几辆车停在路边”；详细描述则需要补充车辆、人物、招牌与相互位置。不同长度并不意味着观察事实不同，但可能使更多属性和关系进入输出，也就增加了需要核验的内容。

\subsubsection{Bottom-Up/Top-Down的区域注意力描述}
\label{sec:vision-understanding-bottomup}

Bottom-Up/Top-Down先用检测器提出对象与显著区域，再根据正在形成的描述选择相关区域。“自底向上”指区域由图像检测产生，“自顶向下”指当前语言状态决定读取哪些区域。其原实现使用在Visual Genome上训练的Faster R-CNN区域特征，描述模型训练期间固定这些图像特征\scite{vision-bottomup2018}

设$K$个候选区域的特征为$\bm v_i$。描述模型由两个LSTM组成：注意力LSTM读取前一时刻语言LSTM的输出、所有区域的平均特征，以及前一个词的嵌入，形成当前状态$\bm h_t^{\mathrm{att}}$。它据此计算
\begin{equation}
  \begin{aligned}
    e_{i,t}&=\bm w^\top\tanh(\bm W_{\mathrm v}\bm v_i
                             +\bm W_{\mathrm h}\bm h_t^{\mathrm{att}}),\\
    \alpha_{i,t}&=\frac{\exp(e_{i,t})}{\sum_{k=1}^{K}\exp(e_{k,t})},
    \qquad
    \hat{\bm v}_t=\sum_{i=1}^{K}\alpha_{i,t}\bm v_i\text{。}
  \end{aligned}
  \label{eq:vision-caption-attention}
\end{equation}
各区域特征与注意力状态先映射到相同维数，$\bm w$再把每个组合压成标量。归一化权重形成当前视觉信息$\hat{\bm v}_t$，它与注意力状态一起输入语言LSTM；后者的输出经线性层和softmax预测下一个词。

对于“一个人走过停着的汽车”，预测人物相关词时可以更多读取人物区域，描述汽车时再读取车辆区域。若两区域的打分为$(\log3,0)$，权重就是$(0.75,0.25)$；这表示当前计算中两种特征的贡献，并不是“第一个词有75\%概率由该区域证明”。

训练以图像及目标描述记录最小化逐词交叉熵，更新两个LSTM、注意力和输出层；原论文也报告了在交叉熵训练后进一步针对CIDEr评分优化的配置。使用时依赖已生成的词前缀逐步选择区域并解码。候选检测遗漏的对象难以通过区域路径被描述，固定词表和训练分布又限制可以表达的内容。注意力图能展示计算读取位置，但不能独立证明每个生成词都正确。

\subsubsection{视觉语言模型的条件描述}
\label{sec:vision-understanding-vlm-caption}

BLIP-2等模型把视觉表示接入已有语言解码器，连接方式见\ref{sec:vision-models-blip2}节。图像描述适配使用“图像、描述条件、目标文字”记录，使解码损失与所要求的输出对应；已有检查点使用时，则应沿用其支持的提示格式和视觉处理。描述指令改变输出选择，并没有增加图像中的事实。

同一张照片要求“一句话概括”时，可以覆盖主体活动；要求“说明右侧两辆车的差别”时，必须读取局部颜色、形状和相对位置。生成后可以将内容拆成“存在两辆车”“右车为蓝色”“左车位于前方”等判断，逐项回查。若某一细节因分辨率不足不可见，应缩小描述范围，或补充局部观测。

流畅度、覆盖和事实性是不同维度。短句可以忠实却遗漏关键内容，长句可以覆盖广却加入不存在的对象。评价时先固定期望描述范围，再判断哪些必要信息被覆盖、哪些输出得不到图像支持。

\subsection{对象与局部区域的描述}
\label{sec:vision-understanding-region-caption}

全局压缩常弱化小物体细节。局部描述可以增加目标区域的有效分辨率，但对象与参照物的关系可能跨越裁剪边界，需要同时保留来源与必要上下文。

\subsubsection{定位裁剪与来源保持的区域描述}
\label{sec:vision-understanding-crop-caption}

一种组合流程先按查询取得区域框，裁剪后记录“源图编号、原图边界、缩放比例”，再将局部图和整图按模型支持的多图接口组织为输入。若模型本来接受框或区域标记，也可以使用其原生接口。训练区域描述器需要对应的“区域—描述”记录，不能把整图描述训练无条件当作已经学过所有区域指代。

假设两辆外观相近的车分别位于左右两侧，局部裁剪显示右车有车顶行李架；整图则显示它位于另一辆车后方。描述“右侧带行李架的车停在另一辆车后面”同时使用了局部外观和全局位置。若只保留右车裁剪，“在后面”的参照就被删掉了。结果应明确当前描述属于哪幅源图的哪个对象，不能把局部图中的左、右重新当作全图关系。

核验包括区域是否选对、局部属性是否读对、引用的关系是否仍有上下文支持。增加裁剪数量会增加读取成本，也可能产生相互矛盾的局部判断，因此应记录裁剪预算，避免只比较最终文字长度。

\subsection{事件与步骤的概括}
\label{sec:vision-understanding-event-caption}

视频概括会舍弃重复和细节，但事件次数与先后可能正是问题需要的信息。压缩的对象和保留条件应由请求决定，不能默认所有重复事件都可合并。

\subsubsection{时间片段描述与事件层次概括}
\label{sec:vision-understanding-hierarchical-caption}

可以先用\ref{chap:vision-video-recognition}章的取样与片段表示，为每段保留实际时间，再交给支持视频或交错帧输入的语言接口形成描述，例如\ref{sec:vision-models-qwen3vl}节的接口。长视频也可以先生成局部描述，再按问题组织总体概括；这是一种分层组合，摘要应仍可追溯到片段。

设教学视频在$[1,2)$秒和$[5,6)$秒各发生一次敲门，$[3,4)$秒有人讲话。问“发生了什么”可以概括为“人物两次敲门，中间说了一句话”；问“第二次敲门是否在讲话之后”则必须保留第二次事件和讲话结束位置。只存储“人物敲门并讲话”会丢掉次数与顺序，即使该短描述本身没有虚构对象。

片段合并应避免把原顺序按主题重新排列，也不能把同一事件在相邻窗口中的重复描述当成两次事件。旁白、字幕和可见动作分别标明来源：字幕出现“我开门了”，不证明镜头中已经呈现开门。保存局部到总体的对应，可以在后续问答时重新取得被概括省略的证据。

\section{视觉问答与证据推理}
\label{sec:vision-understanding-qa}

\term{视觉问答}（visual question answering, VQA）依据图像、视频或文档等视觉内容回答问题。它可能只需读出一个日期，也可能需要连接多个时刻的事件。检索、定位和读取可以成为问答的组成步骤，但最终答案还取决于这些证据如何进入选择、计算或语言生成。

\subsection{问题与答案要求的确定}
\label{sec:vision-understanding-question}

存在问题要求判断是否有目标，属性问题要求给出颜色或状态，计数问题要求定义可数对象和范围，空间问题要求确定参照系，文档问题还需明确字段和单位。对“总额是多少”，应区分读取票据上已经写出的总额，还是依据若干金额自行求和；这两种请求使用的操作不同。

问答模型容易利用问题中的语言规律猜测答案。VQA v2为同一问题收集答案不同的相似图像，使仅凭问题作答的捷径受到限制。对同一问题更换图像后，答案是否随相关证据改变，是有用的诊断，但需要确认新图像确实改变了答案条件\scite{vision-vqav22017}

GQA利用场景图及问题程序构造组合问题，并检查一致性、视觉依据和答案合理性。它能够系统改变属性、关系和组合深度，但这种构造也规定了题目分布，不能代表所有真实用户问题。面向连续对话时，还要保留图像编号、当前指代和已有回答的证据范围，不能只把此前答案作为不容质疑的事实\scite{vision-gqa2019}

\subsection{证据的获取与组织}
\label{sec:vision-understanding-evidence}

获取证据应由当前缺少的事实决定。读数问题需要足够分辨率，事件先后问题需要相应时间片段，多页文档问题则可能先要确定页面。每次增加输入，都应能够说明它补上了哪项信息。

\subsubsection{检索、定位与读取的证据组合}
\label{sec:vision-understanding-evidence-composition}

对“表格中甲的2月金额是多少”，若已给清晰单页，可以直接定位目标单元并读取；若给一份长文档，需要先找到包含该表的页面；若直接提供真实数据表，页面搜索和像素读取已由外部完成。三种条件的最终答案都可能是18元，但视觉工作量不同。

一种无需新增训练的组合规则是：先按标题或内容检索候选页，再读取表格及其行列表头，按“甲、2月”定位单元，最后返回数值、单位和来源。每条中间记录保留源文件或图像编号、页码、框或时间区间、原文字和处理方式。若目标页未被召回，应补充检索；若小数点看不清，应取得更清晰裁剪。重复要求模型“仔细思考”并没有增加这两类缺失证据。

图\ref{fig:vision-qa-evidence}将同一教学表中的差值问题拆成可核验操作。这里的选择与计算是外显组合规则，各模型只承担已经说明的定位、读取或生成职责，不假设它们内部一定沿相同步骤推理。

\begin{figure}[htbp]
  \centering
  % 教学组合规则；不表示某个VLM内部必然具有该操作轨迹。
\begin{tikzpicture}[x=1cm,y=1cm,font=\figurefont\small,
  fact/.style={draw=BookBlue,fill=BookBlue!10,line width=0.9pt,
    minimum width=44mm,minimum height=13mm,align=center},
  operation/.style={draw=BookTeal,fill=BookTeal!10,line width=0.9pt,
    minimum width=50mm,minimum height=11mm,align=center},
  arr/.style={-{Stealth[length=2mm]},BookInk,line width=0.9pt,
    rounded corners=1mm}]
  \node[anchor=west] at (0,5.7) {问题：甲2月比1月多多少？};
  \node[fact] (jan) at (2.2,4.4)
    {甲，1月：12元\\页1，行3，列2};
  \node[fact] (feb) at (7.8,4.4)
    {甲，2月：18元\\页1，行3，列3};
  \node[operation] (subtract) at (5,2.25) {$18\text{元}-12\text{元}=6\text{元}$};
  \coordinate (join) at (5,3.2);
  \draw[BookInk,line width=0.9pt,rounded corners=1mm]
    (jan.south) -- (2.2,3.2) -- (join);
  \draw[BookInk,line width=0.9pt,rounded corners=1mm]
    (feb.south) -- (7.8,3.2) -- (join);
  \draw[arr] (join) -- (subtract.north);
  \node[operation,minimum width=60mm,fill=white] (answer) at (5,0.3)
    {甲2月比1月多6元};
  \draw[arr] (subtract.south) -- (answer.north);
\end{tikzpicture}

  \caption{带来源的表格问答教学流程。两个输入事实都指向同一页、同一行及不同月份单元；差值计算保留操作数与单位。任一读数或行列归属出错，都应返回相应证据检查。}
  \label{fig:vision-qa-evidence}
\end{figure}
\FloatBarrier

\subsubsection{长视频历史与多源证据的组织}
\label{sec:vision-understanding-video-evidence}

视频问题需要的片段可能分散在很长的历史中。可以利用\ref{chap:vision-video-recognition}章的取样、检索和MovieChat式压缩记忆获得候选，再重新读取与问题相关的时刻。压缩历史保存了便于搜索的线索，并不保证保留每次手部动作或每个短暂出现的物体。

声音、字幕和图像应各自保留时间来源。画外有人说“放好了”，画面中的手却仍在移动物体，这可能来自说话者不同、语言含义不同或音画错位。支持音画输入的模型可采用\ref{sec:vision-models-qwen-omni}节的接口，具体区间对应可按\ref{sec:vision-understanding-av-time}节检查，不能只把不同来源文本拼在一起就认定它们同时描述同一事件。

LongVideoBench把视频与字幕共同作为较长上下文，要求利用所指片段的细节回答问题。解释其结果时应分别记录帧数、视频时长、字幕及候选选项，避免把由字幕直接给出的答案全部归于视觉记忆。多选题的正确答案也没有同时验证空间框和精确时间区间\scite{vision-longvideobench2024}

\subsection{依据证据形成答案}
\label{sec:vision-understanding-answering}

证据可以进入不同的输出计算：从固定或动态候选中选一个答案，从文档原文中取连续跨度，选择表格单元并执行操作，或者条件生成文字。这些路线的表达范围与可追踪性不同，应按请求和已有监督选择。

\subsubsection{LoRRA的视觉回答与OCR候选选择}
\label{sec:vision-understanding-lorra}

招牌上的罕见地名可能不在训练答案词表中。TextVQA论文的LoRRA在固定答案集合之外加入当前图像的OCR字符串，使模型可以指向刚刚读到的内容。它分别用问题读取图像特征和OCR文字特征，再融合两路表示，预测固定答案与动态文字候选的得分；外部OCR不与该回答模块共同训练\scite{vision-textvqa2019}

设固定答案数为$n$，当前OCR候选数为$m$，输出为$n+m$个logit。前$n$个位置对应固定答案，后$m$个对应当前图像的OCR字符串；选中第$n+j$个位置就复制第$j$个字符串。原实现还保留按原候选顺序排列的注意力权重，避免只做加权平均后丢失“应复制第几个候选”的线索。

训练使用逐候选二元交叉熵，而不是要求唯一位置的softmax交叉熵。原因是同一答案可能同时出现在固定词表和OCR候选中，两个位置都可以获得正监督。若固定词表没有“青山路”，OCR已正确读出它，模型仍可通过动态位置回答；若OCR仅给出“青山”和“路”两个候选，原LoRRA一次只能复制一个候选，不能把它写成任意多词的自回归复制器。

\subsubsection{DocVQA的OCR序列与抽取式跨度回答}
\label{sec:vision-understanding-docvqa}

文档答案往往已经出现在原文中，连续跨度可以保留原有拼写。原始DocVQA的BERT基线把OCR结果按从上到下、从左到右的约定序列化，将问题与上下文共同编码，再预测答案的开始和结束位置。它利用的是OCR文字，不能因此被描述为BERT直接从像素理解布局\scite{vision-docvqa2021}

设上下文词元表示为$\bm h_j$，两个线性头得到开始概率$p_{\mathrm s}(j)$和结束概率$p_{\mathrm e}(j)$。对训练跨度$(s^*,e^*)$，目标为
\begin{equation}
  L_{\mathrm{span}}=-\log p_{\mathrm s}(s^*)-\log p_{\mathrm e}(e^*)\text{。}
  \label{eq:vision-docvqa-span}
\end{equation}
解码在合法的$s\leq e$及长度约束内选择高分跨度，再由词元到OCR字符串的映射恢复答案。原论文用答案字符串在序列中的首次匹配近似构造训练跨度；重复日期或重复编号可能使这一弱标注指向错误位置。

在教学票据中，序列包含“编号A017日期2026-04-03已付100.00总额128.50”，问题“票据日期是什么”应指向日期对应跨度。若OCR已把03读成08，抽取模块即使选对跨度也会输出错误日期；若双栏顺序错乱，连续跨度又可能混入另一栏文字。要在结果中返回来源，可额外保留OCR词元到框的映射，但这不是原字符串准确率已经自动核验的内容。

\subsubsection{VisionTaPas的表格单元选择与聚合}
\label{sec:vision-understanding-visiontapas}

表格问题中的算术可以通过选择单元和操作来形成答案。ChartQA的VisionTaPas用TaPas编码问题及带行列信息的数据表，用ViT编码图表，再通过跨模态编码器交换信息。视觉分支读取文字表格，文字表格分支也读取视觉特征；末端在文字表格分支预测单元选择和聚合操作。

设被选中的数值为$v_1,\ldots,v_k$。选择“求和”时返回$\sum_i v_i$，选择“计数”时返回所选单元数，选择“无聚合”时直接输出单元内容。原ChartQA还在TaPas的操作头中增加差值、比值以及“是”“否”等答案类别；它没有把模型变成可随意嵌套的程序执行器。

对于“甲的2月金额比1月多多少”，正确选择是甲行的18和12，并选择差值，恢复$18-12=6$元。单元选择指向具体来源，使人能够发现模型是否误用了乙行的10。若请求“甲乙2月之和比1月之和多多少”，则需要$(18+10)-(12+8)=8$元，这超出了单次差值作用于两个原始单元的表达，不能仅凭操作头中有“求和”和“差值”就认定支持该嵌套。

训练也应与操作对应。基础TaPas可以利用答案监督学习单元与聚合；ChartQA为新增差值和比值使用从答案与数据表推导的启发式单元监督。例如已知答案6，在表内寻找差为6的数对，但可能存在多组数值满足同一差值，产生有噪标注。真实表格与恢复表格都可进入该接口，后者的错误会沿单元选择和计算传递，应分别报告。

\subsubsection{VL-T5的图表特征与条件答案生成}
\label{sec:vision-understanding-vlt5}

同一ChartQA设置还使用VL-T5直接生成答案。文字输入由问题与序列化数据表组成，视觉输入来自图表中条形、折线等标记的区域特征；论文使用针对图表适配的Mask R-CNN提取这些特征，而不是把原VL-T5在自然图像中的固定对象特征不加区别地沿用。

编码后，文字解码器根据多模态上下文及已有答案前缀预测后续词元，训练以目标答案交叉熵更新相关参数，使用时逐词或束搜索解码。对于前述差值问题，目标字符串可以是“6”；模型的输出空间不受某个聚合操作集合直接限制，但内部生成“6”的过程也不自动暴露18、12和减法。

开放表达与准确计算是两种性质。生成接口可以输出完整解释，却仍可能在已经给对的数据上算错。判断视觉特征是否有贡献，还应固定数据表条件：若表头与数字足以独立回答，正确生成并不证明模型利用了图表图像；问题引用“蓝色柱”时，才可能需要颜色与系列的视觉对应。

\subsubsection{证据操作的多步组合与解释}
\label{sec:vision-understanding-evidence-operations}

需要嵌套计算或跨图比较时，可以把已核对的事实交给明确的操作规则。例如，两张图分别可见3辆和5辆符合条件的车，先分别保存对象编号和区域，再计算差值$5-3=2$。若问题问的是“新增了几辆”，只比较数量不足以回答，还需要跨图身份对应；增加2辆同时离开1辆，会使数量只增加1辆。

对前述敲门视频，可以将事件列表组织为“敲门1：$[1,2)$；讲话：$[3,4)$；敲门2：$[5,6)$”。由第二次敲门起点5大于讲话终点4，可回答“在讲话结束之后”。若两段存在重叠，则需要改用起点、终点或主要持续时段等更明确的关系，而不能强行给出唯一先后。

这种组合可以复用模型和计算工具，不需要虚构一个新的端到端训练目标。每步保留选取对象、读取事实、操作与结果，解释才有可核验的依据。NExT-QA通过时间关系、动作解释与场景问题考察更复杂的问答；其中“为什么”的标注属于对视频动作的语义解释，不等同于通过干预识别因果效应。仅凭“讲话之后敲门”无法证明讲话导致敲门\scite{vision-nextqa2021}

\subsection{答案的核验与证据补充}
\label{sec:vision-understanding-answer-check}

答案可拆成对象、属性、字符串、数值、单位和关系等判断，分别与来源对照。若输出“总额是128.50美元”，原票据只写人民币符号，则数字正确而单位错误；若正确数值来自另一页的同名表格，来源也不成立。

核验可以改变相关观测后重新作答。例如，把甲2月金额从18改成13，差值应从6变成1；保持图像不变而改问乙行，答案应改为2元。前一种检查对输入变化的响应，后一种检查查询绑定，二者都需要重新确定正确答案，不能只观察模型是否“回答不同”。

缺少证据时回到相应获取步骤；来源冲突时保留冲突并限制结论。多轮对话中，“那一行”“他后来怎样”应绑定到已经明确的单元或对象，新增裁剪、页面或视频帧则登记为新的观测。上一轮回答本身不应替代尚未核对的原始事实。

\section{理解结果的核验与研究方向}
\label{sec:vision-understanding-evaluation}

理解系统可能答对而找错证据，也可能位置正确但读错字符。验收需要同时看任务成果、中间结构和实际输入条件。不同基准提供不同诊断，不能用一个总分替代全部能力判断。

\subsection{任务输出与中间结构的正确性}
\label{sec:vision-understanding-output-metrics}

匹配评价应固定候选单位，报告召回和排序；定位要检查区域、点或区间是否满足查询；读取和结构解析则应拆开字符、行列、阅读关系及连接错误。描述与问答还需核验事实覆盖和答案要求。表\ref{tab:vision-understanding-evaluation}给出沿本章例子的对应。

\begin{table}[htbp]
  \centering
  \small
  \begin{tabularx}{\linewidth}{@{}p{19mm}XX@{}}
    \toprule
    成果 & 直接检查 & 仍需单独检查\\
    \midrule
    检索候选 & 完整正例的召回和名次 & 关系是否由同一对象满足\\
    框或区间 & 与目标范围的吻合 & 指代、身份与空目标\\
    OCR字符串 & 字符编辑距离、整串正确 & 数字含义、单位与来源\\
    表格结构 & 行列、跨度及内容归属 & 单元数值和后续运算\\
    描述或答案 & 内容、数值与请求一致 & 支持它的区域和时刻\\
    \bottomrule
  \end{tabularx}
  \caption{多模态理解的分层核验。每行的直接检查只覆盖一种成果，不能自动证明最后一列的条件。}
  \label{tab:vision-understanding-evaluation}
\end{table}
\FloatBarrier

Winoground固定两句文字的词汇集合，只改变词序，并要求与两幅图正确配对，可用于检查角色和关系绑定\scite{vision-winoground2022}
前述“人搬起箱子”和“箱子挡住人”也体现了相同诊断思路：不能只检测是否出现所有名词。基准论文中对当时模型的观察应限于所测设置，不外推为所有新模型的结论。

\term{字符错误率}（character error rate, CER）用恢复目标字符串所需的替换、删除和插入次数衡量读取误差。若参考共含$N>0$个字符，相应次数为$S,D,I$，则
\begin{equation}
  \mathrm{CER}=\frac{S+D+I}{N}\text{。}
  \label{eq:vision-cer}
\end{equation}
分母使用参考字符数，需规定空格、标点和大小写如何处理；插入很多字符时该值可以超过1。将“128.50”读成“123.50”只有一次替换，CER为$1/6$，但金额字段已经错误，因此业务字段还应报告整串或数值正确性。

POPE以对象是否存在的询问检查一类视觉幻觉，即回答加入了图像中不支持的对象\scite{vision-pope2023}
应同时观察错误肯定和错误否定，避免总回答“没有”获得表面的保守。存在性正确仍不保证颜色、数量、动作、字符串或长段自由描述都正确，这些判断需要各自证据。

VSI-Bench进一步考察视频中的空间信息保存与推断\scite{vision-vsibench2025}
视角改变后，“左侧”可能需要从当前相机坐标转换到场景关系；只记住一句“椅子在桌子左边”可能丢失当时参照。空间问答应保留观察视角、对象对应及测量依据，不能仅从语言解释连贯推断三维关系已经恢复。

\subsection{答案与来源证据的对应}
\label{sec:vision-understanding-grounded-evaluation}

若要求答案、时间和空间都正确，可以把单题成功写成三个条件的交集。令$A_i$表示答案正确，$T_i$表示返回时间区间满足阈值，$B_i$表示目标框满足阈值，则在明确标注协议下，联合成功率为
\begin{equation}
  \frac1n\sum_{i=1}^{n}\bm 1[A_i\land T_i\land B_i]\text{。}
  \label{eq:vision-grounded-success}
\end{equation}
这里$\bm1[\cdot]$为指示函数。这个教学指标的目的在于展示交集关系，不代称某个基准的全部计分规则。100题中若80题答案正确，其中50题时间也正确，这50题中又只有30题空间正确，联合成功率就是30\%，不是把三个环节的边际分数任意相乘。

2026年公开的VideoZeroBench通过不同层级分别收紧答案、时间和空间证据要求，具体体现了答案与依据可能分离的问题\scite{vision-videozerobench2026}
报告时还需说明模型是端到端寻找证据，还是已获得人工给定区间或框；额外给定证据改变了任务难度。该预印本的测试范围不能代表所有长视频应用。

% 完整作者边注不可跨页，为论述开头和来源共同预留阅读范围。
\Needspace{11\baselineskip}
解释长度也不是证据量。More Thinking, Less Seeing?研究了所测多模态推理模型中推理长度与视觉忠实性的关系，并提出相应诊断\scite{vision-morethinking2025}
由此需要检查的是延长输出后，新增判断是否仍有视觉依据；不能概括成“解释越长一定越好”或“一定越差”。若模型多写了五段推理，却始终使用读错的13，更多文字仍不能把正确差值恢复为6。

\Needspace{14\baselineskip}
Video-MME-v2用关联问题的成组评价补充单题准确率，使相互矛盾的回答更难被总分掩盖\scite{vision-videommev22026}
例如，模型既说“第二次敲门在讲话后”，又说“讲话发生在两次敲门之后”，在上述事件记录下至少一项错误。成组一致性是必要诊断，却不能独立证明内容正确：模型也可能一致地记错所有时间。其2026年版本应与所用字幕、音频和帧采样条件一起记录。

\subsection{模态贡献与输入协议}
\label{sec:vision-understanding-input-protocol}

要判断视觉是否贡献了答案，可以移除、替换或错配图像，同时固定问题和其他输入。相同方法也适用于声音、字幕、OCR和数据表。移除模态可能产生训练分布之外的输入，性能变化不宜直接解释为纯粹因果贡献；多种受控替换与具体失败检查能提供更可靠的诊断。

若字幕已经明确写出答案，无图条件下答对并不意外。若真实数据表包含全部行列和数字，系统就不再承担从图像恢复表格的工作。报告需把原始像素、外部OCR、字幕、真实结构以及检索裁剪预算逐项固定，才能比较两个系统是否在相同信息条件下工作。

关键证据丢失、关系绑定错误和多步空间推断相互影响。Video-MSR在2026年将约束定位、链式指代、路线及反事实物理问题组织为多跳空间推理测试，可用于区分单步识别和组合推断。基准内训练带来的改善仍需放回具体任务评估，不能直接推论真实导航可靠性\scite{vision-videomsr2026}

几何证据则提供另一种补充。Geo3R讨论了透视、物体朝向和视角变化引起的空间回答偏差，并在推理时引入几何信息及结构化三维关系\scite{vision-geo3r2026}
作为2026年公开案例，它说明可把空间判断重新落到几何对象上检查；若几何估计本身错误，增加三维表示也可能传播错误，而不是普遍消除视觉幻觉。

这些研究方向共同要求保留可核验的中间对象。读取越多并不总是越有用，关键在于缺少哪项事实、增加的观测是否补上它，以及不同来源冲突时怎样限定答案。持续改进应围绕具体失败设计数据和评价，而不把更长上下文或更多模态本身当作完成理解的标准。

\section{本章小结}
\label{sec:vision-understanding-summary}

从街景定位、票据读取到视频事件先后，多模态理解都需要把请求落到特定观测上。匹配缩小候选范围，定位确定相关区域和时间，读取与解析恢复字符串及关系，答案再由选择、抽取、计算或生成形成。不同接口可以组合，但每次压缩、裁剪和结构转换都可能改变后续可利用的证据。

可检查的结果应同时保留内容和来源。正确数字可能取自错误单元，正确描述可能遗漏关键时序，流畅解释也可能围绕错误观察展开。按任务成果、中间证据和输入协议分层核验，才能确定问题发生在哪里。下一章转向媒体制作：依据给定条件创作新内容时，需要保持的主体、空间和时间关系仍须明确，但未观测细节的生成与本章对既有证据的解释承担不同职责。
